\section{Introduction}

The purpose of these notes is to organize several constructions in equivariant
geometry from the viewpoint of stratified homotopy theory. The guiding principle
is that a geometric object with symmetry should not be studied only as a space,
nor only through its ordinary quotient. Instead, if a group \(G\) acts on a
stratified space \((X,\mathcal S)\), the natural object is the quotient stack
\([X/G]\), equipped with the induced stratification. The corresponding
exit-path \(\infty\)-category then records both the incidence relations among
the strata and the homotopy-coherent information coming from stabilizers and
monodromy.

Thus the basic pattern considered in these notes is
\[
G\curvearrowright (X,\mathcal S)
\quad\Longrightarrow\quad
[X/G]
\quad\Longrightarrow\quad
\operatorname{Exit}([X/G])
\quad\Longrightarrow\quad
\operatorname{Sh}_{\mathrm{cstr},G}(X).
\]
Here \(\operatorname{Sh}_{\mathrm{cstr},G}(X)\) denotes the category of
\(G\)-equivariant constructible sheaves on \(X\), constructible with respect to
the stratification \(\mathcal S\). The point of the exit-path construction is
that constructible sheaves on a stratified space are controlled by functors out
of its exit-path \(\infty\)-category. In the equivariant setting, one should
apply this principle not merely to \(X\), but to the quotient stack \([X/G]\).
This retains the stabilizer data which would be lost by passing to the ordinary
topological quotient.

The affine Grassmannian provides the main example motivating this perspective.
For a complex reductive group \(G\), its affine Grassmannian is
\[
\operatorname{Gr}_G = G((t))/G[[t]].
\]
It carries a natural \(G[[t]]\)-action, and the corresponding orbit
decomposition gives the Schubert stratification
\[
\operatorname{Gr}_G
=
\bigsqcup_{\lambda\in X_*(T)^+}
G[[t]]\cdot t^\lambda .
\]
The geometric categories appearing in the geometric Satake correspondence are
built from sheaves that are constructible, and often equivariant, with respect
to this Schubert stratification. From the present viewpoint, one regards
\(\operatorname{Gr}_G\) as an equivariantly stratified object and studies the
stack
\[
[\operatorname{Gr}_G/G[[t]]],
\]
together with its exit-path \(\infty\)-category.

This framework also clarifies how other examples fit into the same pattern.
In Morse theory, a function or moment map produces a stratification by stable
or unstable manifolds. In Bruhat--Tits theory, buildings are stratified by
facets, alcoves, and apartments, with group actions encoded by parahoric
stabilizers. In toric geometry, the algebraic torus action gives the orbit
stratification, while the compact torus action leads to moment-map and
Morse-type decompositions. These examples are not identical, but they share a
common homotopical skeleton: a space with symmetry, a compatible
stratification, and a category of constructible objects controlled by exit
paths.

The aim of these notes is therefore not to develop the full technical machinery
of equivariant stable homotopy theory, nor to prove the geometric Satake
equivalence. Rather, the goal is to isolate the conceptual mechanism by which
group actions and stratifications combine to produce categories of
constructible sheaves. The affine Grassmannian will serve as the central test
case, while Morse theory, Bruhat--Tits buildings, and toric varieties provide
parallel examples illustrating the same equivariant stratified principle.